%============================================================
% Chapter 11: Gray Codes, Hypercubes and Bit Orders
% STATUS: PLANNED  (generated by tools/outline.py; replace this file when the chapter is written)
%============================================================
\chapter{Gray Codes, Hypercubes and Bit Orders}
\label{ch:gray}

\begin{tcolorbox}[enhanced, colback=white, colframe=secondaryblue, boxrule=1pt, arc=2mm]
\textbf{\large\textcolor{primaryblue}{Order the $2^n$ masks so that neighbours differ in one bit: Gray codes, walks on the hypercube and other useful bit orders.}}
\end{tcolorbox}

\chapterfacts{1500}{2400}{4}{8}{30}

\begin{plannedbox}[Planned Contents --- this chapter is not written yet]
\textbf{Part A --- Lessons}
\begin{enumerate}[leftmargin=*, label=\textbf{A\arabic*.}, itemsep=1pt]
    \item The reflected Gray code $g = x \oplus (x \gg 1)$ and its inverse
    \item The hypercube $Q_n$: neighbours differ in one bit; Hamiltonian cycles from Gray codes
    \item Operations that flip one bit under rules: the Gray-code rank as the answer
    \item Bit-reversal permutation (the FFT order) and other orders of masks
    \item Hamming balls and distances on the hypercube
\end{enumerate}
\textbf{Part B --- Patterns} (with candidate problems, checked when the chapter is written)
\begin{enumerate}[leftmargin=*, label=\textbf{P\arabic*.}, itemsep=1pt]
    \item \textbf{Gray-code ranks} --- Minimum One Bit Operations to Make Integers Zero {\small\color{darkgray}(LeetCode 1611)}
    \item \textbf{Walks on the hypercube} --- Magical Permutation {\small\color{darkgray}(CF 1163E)}
    \item \textbf{Bit-reversal and other orders}
    \item \textbf{Hamming balls}
\end{enumerate}
\end{plannedbox}
